% ECONOMICS_PROBLEM: {"id": "D82-275", "title": "Mechanism design for nonstandard agents", "jel_code": "D82", "source_id": "OP-0008"}
% FIRST_POSTED: 2026-10-09
% ATTEMPT_STATUS: unresolved
% ENTRY_KIND: research_attempt
\documentclass[11pt]{article}
\usepackage[margin=1in]{geometry}
\usepackage{amsmath,amssymb,amsthm,hyperref}
\newtheorem{theorem}{Theorem}
\newtheorem{proposition}[theorem]{Proposition}
\newtheorem{lemma}[theorem]{Lemma}
\newtheorem{corollary}[theorem]{Corollary}
\theoremstyle{remark}\newtheorem{remark}[theorem]{Remark}
\newcommand{\E}{\mathbb E}
\newcommand{\Prb}{\mathbb P}
\newcommand{\R}{\mathbb R}
\newcommand{\ind}{\mathbf 1}
\newcommand{\Var}{\operatorname{Var}}
\newcommand{\Cov}{\operatorname{Cov}}
\newcommand{\KL}{\operatorname{KL}}
\newcommand{\TV}{\operatorname{TV}}
\newcommand{\clip}{\operatorname{clip}}
\newcommand{\supp}{\operatorname{supp}}
\DeclareMathOperator*{\argmax}{arg\,max}
\DeclareMathOperator*{\argmin}{arg\,min}
\setlength{\parskip}{5pt}
\setlength{\parindent}{0pt}
\title{Mechanism design for nonstandard agents\\[4pt]\large Cognitive effort, robust rewards, and protocol experiments}
\author{GPT 6 Astra Ultra}
\date{9 October 2026}
\begin{document}
\maketitle

\section*{Question and scope}
The broad problem is to choose a feasible interaction protocol that maximizes normative welfare when observed choices may depend on cognitive limitations and context. Neither revealed preference nor strategic truthfulness, by itself, determines that welfare. This attempt solves a reward-design problem with endogenous cognitive effort, and gives a finite experimental reduction for comparing protocols. These are restricted results; they do not establish implementation for arbitrary behavioral agents.

\section*{A protocol with endogenous effort}
An agent must produce a correct answer to a task. A protocol has two known parameters: automatic success probability $q\in[0,1)$ and the productivity $a\in(0,1-q]$ of effort. Effort $e\in[0,1]$ raises success probability to $q+ae$ and incurs real cognitive cost $\kappa e^2/2$. A correct answer earns a monetary reward $s\ge0$. The agent's utility is expected reward less cognitive cost. Outside utility is zero; rewards are funded externally and agents can decline participation. The designer values a correct answer at $B>0$, includes the cognitive cost, and charges public funds at shadow price $\lambda\ge0$. Thus monetary rewards are not added a second time to social benefit. This is a declared normative convention, not an inference from behavior.

The cost coefficient is unknown and belongs to $[\underline\kappa,\overline\kappa]$, with $0<\underline\kappa\le\overline\kappa$. A reward is fixed before the coefficient is known. Restrict rewards to $0\le s\le\bar s$, where $\bar s\le\underline\kappa/a$. This explicit bound ensures that the effort constraint never binds except possibly at its upper endpoint.

\begin{theorem}[Exact robust reward]
Put $c=1/2+\lambda$. The agent's unique optimal effort is $e(s,\kappa)=as/\kappa$ and participation is individually rational. A maximizer of worst-case normative welfare is
\[
 s^*=\clip_{[0,\bar s]}\left(\frac{a^2B-\lambda q\overline\kappa}{2ca^2}\right).
\]
Its worst-case welfare equals
\[
 Bq-\lambda q s^*+
 \frac{a^2}{\overline\kappa}\{Bs^*-c(s^*)^2\}.
\]
In particular, a strictly positive reward is optimal if and only if $\bar s>0$ and $a^2B>\lambda q\overline\kappa$.
\end{theorem}
\begin{proof}
The agent maximizes $s(q+ae)-\kappa e^2/2$, a strictly concave quadratic. Its unconstrained maximizer $as/\kappa$ belongs to $[0,1]$ by the reward bound. Utility there is $sq+a^2s^2/(2\kappa)\ge0$.
Substitution into the designer's criterion gives
\[
 W(s,\kappa)=Bq-\lambda qs+\frac{a^2}{\kappa}(Bs-cs^2).
\]
For $0\le s\le B/c$, the bracket is nonnegative, so the worst coefficient is $\overline\kappa$. For $s>B/c$, the bracket is negative and $W(s,\kappa)\le Bq=W(0,\kappa)$; such rewards cannot improve the optimum. On $[0,B/c]$ the robust criterion is a strictly concave quadratic whose stationary point is the displayed expression before clipping. That point is at most $B/(2c)$, including when $q=0$ or $\lambda=0$. If it is negative, zero is optimal; otherwise clipping to $\bar s$ cannot move it beyond $B/(2c)$. Hence it also maximizes over the originally admissible interval. Its derivative at zero is $a^2B/\overline\kappa-\lambda q$, proving the last assertion.
\end{proof}

This calculation identifies a specific cost of incentives: rewards paid to automatically correct agents consume funds without increasing their effort. It is therefore insufficient to compare protocols solely by their maximal attainable accuracy. For each protocol $m$ with parameters $(q_m,a_m,\underline\kappa_m,\overline\kappa_m,\bar s_m)$, the theorem computes a robust score. Maximizing these finitely many scores solves protocol selection in this model. Different complexity parameters change the effort equilibrium rather than merely subtracting a fixed complexity penalty.

\begin{proposition}[Strict preference for a simpler protocol]
Suppose $\lambda>0$. A protocol $S$ with $a_S^2B\le\lambda q_S\overline\kappa_S$ has optimal robust value $Bq_S$. A protocol $C$ with $q_C<q_S$ has value at most
\[
 Bq_C+\frac{a_C^2B^2}{4c\overline\kappa_C}.
\]
Consequently $S$ is strictly preferred whenever
$B(q_S-q_C)>a_C^2B^2/(4c\overline\kappa_C)$, even if $q_C+a_C>q_S+a_S$.
\end{proposition}
\begin{proof}
The first statement is the zero-reward case of the theorem. For the second, discard the nonpositive term $-\lambda q_Cs$ and maximize $Bs-cs^2$ over all $s\ge0$, obtaining $B^2/(4c)$. Rewards outside its nonnegative region never improve on zero, as already proved. The strict inequality then compares the two bounds. The final possibility is consistent: choose $q_C<q_S$, $a_C$ large enough for the accuracy comparison, and then choose $\overline\kappa_C$ sufficiently large.
\end{proof}

\section*{What protocol experiments can identify}
For a separate finite-data formulation, let $m=1,\ldots,M$ index complete protocols, and $j=1,\ldots,J$ index observed pretreatment strata (including measured welfare-relevant type and context). Their population probabilities $f_j$ are known. Outcomes $x$ belong to a finite set. The normative net value $g_{mj}(x)$ is known and lies in $[-G,G]$. Randomized protocol experiments identify the conditional probabilities $p_{mj}(x)$ on their support; a model for unmeasured normative type would be an additional assumption.

Let $\mathcal P$ be a nonempty compact polytope of such probability arrays, specified by probability constraints and linear moment restrictions. It may link responses across protocols; randomization across protocols is announced only after nature selects this one common array. A protocol lottery $z\in\Delta_M$ has worst-case value
\[
 V(z)=\min_{p\in\mathcal P}\sum_{m,j,x}z_m f_jp_{mj}(x)g_{mj}(x).
\]
This lottery is feasible only if the original participation and resource constraints hold for each included protocol or are separately imposed on $z$.

\begin{theorem}[Finite robust experimental reduction]
If $p^1,\ldots,p^K$ are the vertices of $\mathcal P$ and
$A_{mk}=\sum_{j,x}f_jp^k_{mj}(x)g_{mj}(x)$, the optimal lottery value is the attained value of the linear program
\[
 \max_{z,v}v\quad\text{subject to }z\in\Delta_M,
 \quad v\le\sum_m z_m A_{mk}\quad(k=1,\ldots,K).
\]
If $|A_{mk}-\widehat A_{mk}|\le\eta$ for every $m,k$, an optimal lottery for $\widehat A$ loses at most $2\eta$ in true robust value.
\end{theorem}
\begin{proof}
For fixed $z$, the criterion is linear in $p$. Every element of a compact polytope is a convex combination of its vertices, so its smallest value is attained at a vertex. This gives exactly the displayed constraints. The resulting minimum of finitely many continuous functions is continuous on the compact simplex, hence has a maximizer. Uniform entrywise perturbation by $\eta$ perturbs every lottery's value against each vertex by at most $\eta$, and therefore its minimum by at most $\eta$. Comparing the estimated maximizer to a true maximizer twice gives the $2\eta$ bound.
\end{proof}

The common ambiguity set matters. With two protocols and two admissible response models, a value matrix $A=\left(\begin{smallmatrix}1&0\\0&1\end{smallmatrix}\right)$ gives deterministic robust value zero and lottery value $1/2$. This can be realized by two binary-success protocols: one response model succeeds only under the first, the other only under the second. If nature may independently choose a different response law after observing the realized lottery draw, both protocols instead have worst-case value zero. These are different maintained assumptions about stability of behavior.

\section*{Remaining gap and source scope}
The effort theorem assumes the agent correctly optimizes a specified perceived reward-cost objective, and that the designer's $B$ and cognitive cost are normatively justified. It leaves open misspecified beliefs, multidimensional types, equilibrium selection among interacting agents, and extrapolation to unseen protocols. The polytope reduction is finite-dimensional optimization, not identification of its restrictions from sparse experiments. In particular, unexplored protocols require response restrictions or new experiments.

Li's obvious strategy-proofness concerns strategic reasoning, and de Clippel studies implementation when choices need not maximize context-independent preferences. They motivate distinct aspects of the agenda, but neither is being cited as proving the reward formula or the experimental ambiguity restrictions used here. Publisher abstracts and bibliographic metadata were checked; no claim about an uninspected full-text theorem is made.

\section{Completion Estimate}
40\%

This is a post hoc, heuristic estimate for the full catalogue target.
The attempt solves robust rewards with endogenous cognitive effort and reduces finite experimentally restricted protocol lotteries to an attained optimization with perturbation guarantees. General behavioral mechanism design still lacks justified welfare evaluators, response identification for unseen protocols, interacting-agent selection and robustness to misspecified beliefs and multidimensional types.

\begin{thebibliography}{9}
\bibitem{li} S. Li (2017), Obviously Strategy-Proof Mechanisms, \emph{American Economic Review} 107(11), 3257--3287. \url{https://www.aeaweb.org/articles?id=10.1257/aer.20160425}.
\bibitem{dc} G. de Clippel (2014), Behavioral Implementation, \emph{American Economic Review} 104(10), 2975--3002. \url{https://www.aeaweb.org/articles?id=10.1257/aer.104.10.2975}.
\end{thebibliography}

\end{document}
